\documentclass[11pt]{article}

\usepackage[T1]{fontenc}
\usepackage{mathtools,amsthm,amssymb}
\usepackage{lmodern}
\usepackage{microtype}
\usepackage[margin=1.05in]{geometry}
\usepackage{enumitem}
\usepackage{needspace}
\usepackage{xcolor}
\usepackage{tikz}
\usetikzlibrary{arrows.meta}
\pdftrailerid{}
\pdfsuppressptexinfo=15
\usepackage{hyperref}
\hypersetup{colorlinks=true,linkcolor=black,citecolor=black,
  urlcolor=blue!45!black,pdftitle={The second Kahn--Kalai conjecture},pdfauthor={OpenAI}}

\newtheorem{theorem}{Theorem}[section]
\newtheorem{lemma}[theorem]{Lemma}
\numberwithin{equation}{section}

\newcommand{\PP}{\mathbb P}
\newcommand{\EE}{\mathbb E}
\newcommand{\pc}{p_{\mathrm c}}
\newcommand{\pE}{p_{\mathrm E}}
\newcommand{\ind}{\mathbf 1}
\newcommand{\TT}{\mathcal T}
\newcommand{\UU}{\mathcal U}
\newcommand{\ch}{\operatorname{Ch}}
\newcommand{\ee}{\mathrm e}
\newcommand{\unverified}[1]{\par\smallskip\noindent\fbox{\parbox{0.95\linewidth}{\textbf{Unverified step.} #1}}\par\smallskip}
\newcommand{\step}[1]{\par\smallskip\noindent\textbf{Step #1.}\ }

\title{The second Kahn--Kalai conjecture}
\author{Rewritten for readability from the OpenAI manuscript of the same title}
\date{September 24, 2026}

\begin{document}
\maketitle

\section{The result}\label{sec:introduction}

Graphs are finite and simple. $N(G,F)$ is the number of subgraphs of $G$
isomorphic to $F$ (ordinary copies: $G$ may have further edges among the
chosen vertices). $G(n,p)$ is the random graph on $[n]=\{1,\ldots,n\}$ in
which each edge is present independently with probability $p$. For a graph
$H$ with at least one edge and at most $n$ vertices, define
\begin{align*}
 \pc(n,H)&=\inf\{p\in[0,1]:\PP(N(G(n,p),H)>0)\ge\tfrac12\},\\
 \pE(n,H)&=\inf\{p\in[0,1]:\EE N(G(n,p),F)\ge\tfrac12
                 \text{ for every subgraph }F\subseteq H\},
\end{align*}
the \emph{containment threshold} and the \emph{graph expectation
threshold}. Three elementary facts:
\begin{enumerate}[label=(\alph*),itemsep=1pt,topsep=3pt]
 \item\label{fact:a} $\EE N(G(n,p),F)=N(K_n,F)\,p^{e(F)}$, since each copy
 of $F$ in $K_n$ is present with probability $p^{e(F)}$. This expectation
 and the probability defining $\pc$ are polynomials in $p$, and every
 constraint holds at $p=1$; hence both infima are attained.
 \item $\pE(n,H)>0$: the single-edge subgraph requires
 $\binom n2p\ge\tfrac12$.
 \item $\pE(n,H)\le\pc(n,H)$: a graph containing $H$ contains every
 $F\subseteq H$, so at $p=\pc(n,H)$ every $F$ satisfies
 $\EE N(G(n,p),F)\ge\PP(N(G(n,p),F)>0)\ge\tfrac12$.
\end{enumerate}

\begin{theorem}\label{thm:main}
For every $n\ge2$ and every finite simple graph $H$ with $h=e(H)\ge1$
edges and at most $n$ vertices,
\begin{equation}\label{eq:stronger}
 \pc(n,H)\le
 \min\{1,\;2048\ee^{50}\pE(n,H)(1+\log_2 h)\}.
\end{equation}
In particular,
$\pc(n,H)\le6144\ee^{50}\pE(n,H)\log_2 n$.
\end{theorem}

The constants are universal and $H$ may vary with $n$. Kahn and
Kalai~\cite{KK07} use expected count $1$ instead of $\tfrac12$;
by~\ref{fact:a} this multiplies the critical density of a subgraph with
$x\ge1$ edges by $2^{1/x}\le2$, so it changes $\pE$ by a factor in $[1,2]$
and gives the same conjecture.

\paragraph{Conventions.}
$\ee$ is the base of the natural logarithm $\log$. For a finite set $X$
and $\rho\in[0,1]$, $\mu_\rho$ is the law of the random subset of $X$
containing each element independently with probability $\rho$.
\emph{Monotonicity} means: if a family of subsets of $X$ is closed under
supersets, its $\mu_\rho$-probability is nondecreasing in $\rho$ (couple
the two samples so that one contains the other).
To keep one meaning per symbol we renamed some of the source's letters:
$K,i$ (source: $R,r$) for intersection sizes in Lemma~\ref{lem:local},
$\tau$ (source: $t$) for the number of successes in
Section~\ref{sec:trees}, $x$ (source: $r$) for $|S|$ in
Section~\ref{sec:graphs}, and $\kappa$ for the transition probabilities,
which the source leaves unnamed.

\subsection*{Proof at a glance}

\emph{Key idea.} A copy of $H$ is grown in stages along a chain of edge
sets $\varnothing=H_0\subset\cdots\subset H_k=H$. All growth histories
form a tree; each arc is labelled by the edges added at that stage. One
shows that a random edge set contains all labels along some root-to-leaf
path. The same random set is reused at all stages, which is affordable
because labels along one history are disjoint.

\begin{enumerate}[label=\arabic*.,itemsep=2pt,topsep=3pt]
 \item \textbf{Lemma~\ref{lem:local} (one family).} Call a random set $A$
 \emph{$a$-spread} if $\PP(J\subseteq A)\le a^{|J|}$ for all nonempty $J$.
 Given an $a$-spread family of sets of size at most $m$ and
 $W\sim\mu_\rho$ with $\rho\ge\ee^{50}a$, a randomized rule selects a
 member $A_b$ whose \emph{fragment} $A_b\setminus W$ has size at most
 $m/2$, except with probability $\ee^{-9m}$. A change-of-measure bound
 says that selecting $b$ distorts the law of $W$ outside $A_b$ by at most
 a factor $\ee^{11m}$.
 \item \textbf{Lemma~\ref{lem:reduction} (one reduction)} applies
 Lemma~\ref{lem:local} at every node of a tree with one common $W$. With
 probability at least $3/4$ it returns a subtree labelled by fragments:
 all label sizes are halved, spread worsens by a factor close to $1$, and
 $W$ together with a path of the new tree covers a path of the old one.
 Disjointness along paths is what makes the change-of-measure bound
 applicable down the tree.
 \item \textbf{Theorem~\ref{thm:tree} (tree covering)} iterates
 Lemma~\ref{lem:reduction} with fresh sets. After
 $\lfloor\log_2\ell_k\rfloor+1$ successes all labels are empty, so the
 union of the sets used contains a path of the original tree. This is the
 origin of the factor $1+\log_2h$.
 \item \textbf{Lemma~\ref{lem:hierarchy} (graph hierarchy)}, independent
 of items 1--3, builds the chain and the tree for $H$: every stage is
 $128\pE$-spread and label sizes grow by a factor at least $16$.
 \item \textbf{Theorem~\ref{thm:main}} is Theorem~\ref{thm:tree} applied
 to the tree of Lemma~\ref{lem:hierarchy}; note $16\cdot128=2048$.
\end{enumerate}
Dependencies: Lemma~\ref{lem:local} $\Rightarrow$
Lemma~\ref{lem:reduction} $\Rightarrow$ Theorem~\ref{thm:tree};
Theorem~\ref{thm:tree} and Lemma~\ref{lem:hierarchy} $\Rightarrow$
Theorem~\ref{thm:main}.

\subsection*{Context (quoted from the source; not used in the proof)}

Thresholds for fixed targets are classical~\cite{ER60,Bollobas81}; here
$H$ may grow with $n$. For a perfect matching on $n$ vertices,
$\pE=\Theta(1/n)$ but $\pc=\Theta(\log n/n)$~\cite{ER66,KK07}, and
Hamilton cycles show the same gap~\cite{KK07}, so a logarithm is needed;
other spanning targets, such as the cube~\cite{Riordan00}, have required
different methods.

For an increasing family $\mathcal F\subseteq2^X$, a \emph{cover} is a
family $\mathcal G$ such that every member of $\mathcal F$ contains a
member of $\mathcal G$; its cost at $p$ is $\sum_{S\in\mathcal G}p^{|S|}$,
and $q(\mathcal F)$ is the largest $p$ at which some cover costs at most
$\tfrac12$. Fractional covers define $q_f(\mathcal F)\ge q(\mathcal F)$
similarly~\cite[Section~6]{Tal10}. The bound
$p_{\mathrm c}(\mathcal F)=O(q_f(\mathcal F)\log\ell)$ was proved
in~\cite{FKNP21}, using the fragment method of~\cite{ALWZ21,LSZ19}, and
the same bound with $q$ (the abstract Kahn--Kalai conjecture)
in~\cite{PP24}; $q$ and $q_f$ agree up to a constant factor
by~\cite[Theorem~1.1]{OpenAIThresholdEquivalence2026}. For the family
$\mathcal F_H$ of edge sets containing a copy of $H$, these results alone
do not give Theorem~\ref{thm:main}, because
$\pE(n,H)\le q(\mathcal F_H)$. Reason: by~\ref{fact:a}, $\pE$ is the
largest, over subgraphs $F\subseteq H$ with an edge, of the density at
which $\EE N(G(n,p),F)=\tfrac12$; the copies of a maximizing $F$ form a
cover of $\mathcal F_H$ whose cost at $p$ is $\EE N(G(n,p),F)$, which
equals $\tfrac12$ at $p=\pE$.

Mossel, Niles-Weed, Sun and Zadik~\cite[Theorem~1]{MNSZ22} proved a
single-logarithm bound relative to a modified threshold (expected count
at least $N(H,F)/2$), which equals
$q_f(\mathcal F_H)$~\cite[Equation~(1)]{DKP25}. Dubroff, Kahn and Park
proved $\pc=O(\pE\log^3n)$, and the single logarithm when
$\pE<1/(3n)$~\cite[Theorems~1.2--1.5]{DKP25}; they also compared $N(H,F)$
with expected counts when $F$ is a clique, a cycle or a bounded-degree
tree~\cite[Theorems~1.4--1.5]{DKP25Cases}. Tran~\cite{Tran26} obtained
$q_f(\mathcal F_H)=O(\pE\log(2e(H)))$, hence a loss $O(\log^2(2e(H)))$ in
general, proved the conjecture for trees and for graphs whose average
degree is at least the logarithm of their maximum degree, and proposed a
coupling conjecture~\cite[Section~4]{Tran26} under which the logarithmic
sampling cost would be paid only once.

\emph{What is new.} Lemma~\ref{lem:hierarchy} follows Tran's
maximization argument~\cite[Section~3]{Tran26} (antecedents:
\cite[Claim~29]{LSZ19}, \cite[Lemma~2.4]{ALWZ21}); in the form proved
below it gives spread $128\pE$ with no logarithm. The transition of
Section~\ref{sec:resampling} is from~\cite{MNSZ25}. The new step is
Theorem~\ref{thm:tree}: all stages are reduced with the same random sets,
so the logarithm is paid once for the whole chain.

\section{Resampling one spread family}\label{sec:resampling}

Throughout this section $X$ is a finite set and $(A_b)_{b\in B}$ is a
finite indexed family of subsets of $X$ (repetitions and empty sets
allowed) with weights $\nu_b>0$, $\sum_b\nu_b=1$. For $a>0$ the family is
\emph{$a$-spread} if
\begin{equation}\label{eq:spread}
 \PP_{c\sim\nu}(J\subseteq A_c)\le a^{|J|}
 \qquad(\varnothing\ne J\subseteq X).
\end{equation}
Fix $0<\rho<1$. For $Y\subseteq X$ and $b\in B$ put
\begin{equation}\label{eq:posterior}
 Z(Y)=\sum_{d\in B}\nu_d\rho^{-|A_d|}\ind_{\{A_d\subseteq Y\}},
 \qquad
 \kappa(b\mid Y)=\frac{\nu_b\rho^{-|A_b|}\ind_{\{A_b\subseteq Y\}}}{Z(Y)}.
\end{equation}

\paragraph{The transition.}
This is the planted-set and posterior-resampling construction
of~\cite[Sections~2.1--2.3]{MNSZ25}. Given a set $W\subseteq X$:
\begin{enumerate}[label=\arabic*.,itemsep=1pt,topsep=3pt]
 \item draw a \emph{source} $c$ with law $\nu$ and set $Y=W\cup A_c$;
 \item given $(W,c)$, draw a \emph{target} $b$ with probability
 $\kappa(b\mid Y)$;
 \item output the \emph{fragment} $T=A_b\setminus W$.
\end{enumerate}
Step~2 makes sense because $Z(Y)\ge\nu_c\rho^{-|A_c|}>0$. Thus, for a
fixed $W$,
\begin{equation}\label{eq:joint}
 \PP(\mathrm{source}=c,\ \mathrm{target}=b\mid W)
 =\nu_c\,\kappa(b\mid W\cup A_c).
\end{equation}
For every $W$ the source has law $\nu$, and so, like the source, does the target.
Three observations:
\begin{enumerate}[label=(O\arabic*),itemsep=2pt,topsep=3pt]
 \item\label{obs:incl} \emph{Inclusions.} If the pair $(c,b)$ has
 positive probability, then $A_b\subseteq W\cup A_c$, hence
 \begin{equation}\label{eq:two-inclusions}
  T\subseteq A_c\cap A_b,\qquad A_b\subseteq W\cup T.
 \end{equation}
 The first inclusion compares the fragment with the source. The second
 says that $W$ and the fragment together cover the target label.
 \item\label{obs:local} \emph{Locality.} Since
 $\ind_{\{A_d\subseteq W\cup A_c\}}=\ind_{\{A_d\setminus A_c\subseteq W\}}$,
 the probabilities \eqref{eq:joint}, the value $Z(Y)$ and the fragment
 depend on $W$ only through $W\cap\bigcup_{d\in B}A_d$.
 \item\label{obs:law} \emph{Law of $Y$.} Now let $W\sim\mu_\rho$,
 independently of $c$. For fixed $c$ and $y\supseteq A_c$, we have
 $W\cup A_c=y$ exactly when $W$ contains $y\setminus A_c$ and misses
 $X\setminus y$ (elements of $A_c$ are unconstrained), so
 \[
  \PP(W\cup A_c=y)=\rho^{|y|-|A_c|}(1-\rho)^{|X\setminus y|}
  =\mu_\rho(y)\rho^{-|A_c|};
 \]
 the probability is $0$ if $A_c\not\subseteq y$. Averaging over
 $c\sim\nu$,
 \begin{equation}\label{eq:density}
  \PP(Y=y)=\mu_\rho(y)Z(y)\qquad(y\subseteq X).
 \end{equation}
 (By Bayes' formula, $\kappa(\cdot\mid Y)$ is then the conditional law of
 the source given $Y$, so the target is a conditionally independent
 resample of the source. This motivates the rule but is not used below.)
\end{enumerate}

\begin{lemma}[Local resampling]\label{lem:local}
Suppose $|A_b|\le m$ for all $b$, where $m$ is a positive integer, the
family is $a$-spread, and
\[
 0<\rho<1,\qquad a/\rho\le\ee^{-50}.
\]
Let $W\sim\mu_\rho$ be independent of the source and perform the
transition; probabilities and expectations below are over $W$, source and
target jointly. Call the outcome a \emph{local failure} if
$Z(Y)<\ee^{-10m}$ or $|T|>m/2$. Then
\begin{equation}\label{eq:local-failure}
 \PP(\text{local failure})\le\ee^{-9m}.
\end{equation}
Moreover, for each $b\in B$ and every function $f\ge0$ of $W$ that depends
only on $W\setminus A_b$,
\begin{equation}\label{eq:distortion}
 \EE\bigl[\ind_{\{\mathrm{target}=b,\ Z(Y)\ge\ee^{-10m}\}}f(W)\bigr]
 \le\ee^{11m}\nu_b\,\EE f(W).
\end{equation}
\end{lemma}

The lemma is used at every node of the tree in Lemma~\ref{lem:reduction}:
\eqref{eq:local-failure} bounds the failure mass, and
\eqref{eq:distortion} is applied with $f$ a potential of the subtree
below~$b$.

\begin{proof}
\emph{Plan.} Small $Z$ is rare by \eqref{eq:density}. When $Z$ is not
small, $\kappa$ is at most $\ee^{10m}\nu_b\rho^{-|A_b|}$ times an
indicator, and everything reduces to the intersection sizes
$|A_c\cap A_b|$, which spread controls.

\step{1 (small denominator)} By \eqref{eq:density},
\[
 \PP\bigl(Z(Y)<\ee^{-10m}\bigr)
 =\sum_{y}\mu_\rho(y)Z(y)\ind_{\{Z(y)<\ee^{-10m}\}}\le\ee^{-10m}.
\]

\step{2 (intersections)} Fix $b$ and let $K=|A_c\cap A_b|$ with
$c\sim\nu$. If $K=i$, then $J=A_c\cap A_b$ is an $i$-element subset of
$A_b$ with $J\subseteq A_c$. There are at most $\binom mi$ such $J$, each
with probability at most $a^i$ by \eqref{eq:spread} (trivially if $i=0$).
Hence, for every set $\mathcal I\subseteq\{0,\ldots,m\}$,
\begin{equation}\label{eq:intersection}
 \sum_c\nu_c\rho^{-|A_c\cap A_b|}\ind_{\{|A_c\cap A_b|\in\mathcal I\}}
 =\sum_{i\in\mathcal I}\PP(K=i)\rho^{-i}
 \le\sum_{i\in\mathcal I}\binom mi(a/\rho)^i.
\end{equation}

\step{3 (cancellation)} For fixed $c,b$, the event
$A_b\subseteq W\cup A_c$ is the event $A_b\setminus A_c\subseteq W$, of
probability $\rho^{|A_b\setminus A_c|}$. Therefore
\begin{equation}\label{eq:cancel}
 \rho^{-|A_b|}\,\PP_W(A_b\subseteq W\cup A_c)=\rho^{-|A_c\cap A_b|}.
\end{equation}

\step{4 (large fragment)} On $\{Z(Y)\ge\ee^{-10m}\}$ we have
$\kappa(b\mid Y)\le\ee^{10m}\nu_b\rho^{-|A_b|}\ind_{\{A_b\subseteq Y\}}$,
and by \eqref{eq:two-inclusions} $|T|>m/2$ forces $|A_c\cap A_b|>m/2$.
Summing \eqref{eq:joint} over pairs and averaging over $W$,
\begin{align*}
 \PP\bigl(Z(Y)\ge\ee^{-10m},\ |T|>m/2\bigr)
 &\le\ee^{10m}\sum_{c,b}\nu_c\nu_b\rho^{-|A_b|}
   \PP_W(A_b\subseteq W\cup A_c)\ind_{\{|A_c\cap A_b|>m/2\}}\\
 &=\ee^{10m}\sum_b\nu_b\sum_c\nu_c\rho^{-|A_c\cap A_b|}
   \ind_{\{|A_c\cap A_b|>m/2\}}
   &&\text{by \eqref{eq:cancel}}\\
 &\le\ee^{10m}\sum_{i>m/2}\binom mi(a/\rho)^i
   &&\text{by \eqref{eq:intersection}}\\
 &\le\ee^{10m}2^m\ee^{-25m},
\end{align*}
since $(a/\rho)^i\le\ee^{-50i}\le\ee^{-25m}$ for $i>m/2$ and
$\sum_i\binom mi=2^m$.

\step{5 (total)} Steps 1 and 4 give
$\PP(\text{local failure})\le\ee^{-10m}+\ee^{-(15-\log2)m}$. Dividing by
$\ee^{-9m}$, this is at most $\ee^{-m}+\ee^{-(6-\log2)m}\le
\ee^{-1}+\ee^{-5}<1$ because $m\ge1$. This proves
\eqref{eq:local-failure}.

\step{6 (change of measure)} Bounding $\kappa$ as in Step~4,
\[
 \EE\bigl[\ind_{\{\mathrm{target}=b,\ Z(Y)\ge\ee^{-10m}\}}f(W)\bigr]
 \le\ee^{10m}\nu_b\sum_c\nu_c\rho^{-|A_b|}
   \EE\bigl[\ind_{\{A_b\setminus A_c\subseteq W\}}f(W)\bigr].
\]
The indicator depends only on $W\cap A_b$ and $f(W)$ only on
$W\setminus A_b$; under the product measure $\mu_\rho$ these are
independent, so the last expectation equals
$\rho^{|A_b\setminus A_c|}\EE f(W)$. By the arithmetic of
\eqref{eq:cancel} and then \eqref{eq:intersection} with
$\mathcal I=\{0,\ldots,m\}$, the right-hand side is at most
\[
 \ee^{10m}\nu_b\,\EE f(W)\sum_c\nu_c\rho^{-|A_c\cap A_b|}
 \le\ee^{10m}\nu_b\,\EE f(W)\,(1+a/\rho)^m
 \le\ee^{11m}\nu_b\,\EE f(W),
\]
using $1+a/\rho\le\ee$.
\end{proof}

\section{Covering a tree}\label{sec:trees}

\paragraph{Definitions.}
A \emph{probability tree on $X$ with $k\ge0$ levels} is a finite rooted
tree whose leaves all have depth $k$ (the root has depth $0$), together
with, for each non-leaf node $v$, a probability distribution $\nu_v$ with
positive weights on its set of children $\ch(v)$, and a label
$A_v(b)\subseteq X$ on the arc from $v$ to each child $b$. \emph{Level
$j$} ($1\le j\le k$) is the set of arcs from depth $j-1$ to depth $j$; a
node at depth $j-1$ has \emph{outgoing level} $j$. Different arcs may
carry equal labels; nodes are never identified. A \emph{path} $P$ always
runs from the root to a leaf, and
$\UU(P)=\bigcup_{(v,b)\in P}A_v(b)$. If $k=0$ the only path is empty, with
$\UU=\varnothing$.

The tree is \emph{disjoint} if the labels along each path are pairwise
disjoint. Level $j$ has \emph{capacity} $m$ if all its labels have size at
most $m$, and is \emph{$a$-spread} if, for every node $v$ with outgoing
level $j$, the family $(A_v(b))_{b\in\ch(v)}$ with weights $\nu_v$ is
$a$-spread in the sense of \eqref{eq:spread}. Nothing is assumed about the
law of a whole path.

\begin{theorem}[Tree covering]\label{thm:tree}
Let $\TT$ be a disjoint probability tree on a finite set $X$, with
$k\ge1$ levels. Suppose that every level is $\sigma$-spread, where
$\sigma>0$, and that the levels have positive integer capacities
$\ell_1,\ldots,\ell_k$ with
\[
 \ell_{i+1}\ge16\ell_i\qquad(1\le i<k).
\]
Set $r=\min\{1,\;16\ee^{50}\sigma(1+\log_2\ell_k)\}$. Then
\[
 \PP_{W\sim\mu_r}
 \bigl(\UU(P)\subseteq W\text{ for some path }P\text{ of }\TT\bigr)
 \ge\tfrac23.
\]
\end{theorem}

It is used once, in Section~\ref{sec:graphs}, with $\sigma=128\pE(n,H)$
and $\ell_k\le e(H)$.

\subsection{One simultaneous reduction}

\begin{lemma}[Reduction]\label{lem:reduction}
Let $\TT$ be a disjoint probability tree with $k\ge1$ levels and positive
integer capacities $m_1,\ldots,m_k$ satisfying $m_{j+1}\ge16m_j$. Suppose
level $j$ is $\alpha_j$-spread, where
\[
 0<\alpha_j\le a,\qquad\rho=\ee^{50}a<1.
\]
There is a deterministic procedure which, given $W\sim\mu_\rho$, succeeds
with probability at least $3/4$ and on success returns a disjoint
probability tree $\TT'$ with the same number $k$ of levels, such that:
\begin{enumerate}[label=(\roman*),leftmargin=2.2em,itemsep=2pt,topsep=3pt]
 \item level $j$ has capacity $\lfloor m_j/2\rfloor$ (possibly $0$) and
 is $\alpha_j/(1-\ee^{-m_j})$-spread;
 \item for every path $Q$ of $\TT'$ there is a path $P$ of $\TT$ with
 \begin{equation}\label{eq:recovery}
  \UU(P)\subseteq W\cup\UU(Q).
 \end{equation}
\end{enumerate}
\end{lemma}

It is the single step iterated in the proof of Theorem~\ref{thm:tree}.

\begin{proof}
\emph{Plan.} Run the transition of Section~\ref{sec:resampling} at every
node with the same $W$. Discard nodes at which too much probability is
lost (Steps 0--1), bound the chance that the root is discarded via a
potential (Steps 2--5), and read off the new tree (Steps 6--7).

\step{0 (local quantities)} Let $v$ be a node with outgoing level $j$.
Apply the transition to the family $(A_v(b))_{b\in\ch(v)}$ with weights
$\nu_v$ and $m=m_j$. Since $\alpha_j\le a$ this family is $a$-spread, and
$a/\rho=\ee^{-50}$, so Lemma~\ref{lem:local} applies. For a fixed $W$,
with probabilities over source and target at $v$, set
\[
 d_v(W)=\PP(\text{local failure}\mid W),\qquad
 t_{vb}(W)=\PP(\text{no local failure, target}=b\mid W).
\]
Then
\begin{equation}\label{eq:local-masses}
 d_v(W)+\sum_{b\in\ch(v)}t_{vb}(W)=1,
\end{equation}
and by~\ref{obs:local} both $d_v$ and $t_{vb}$ depend on $W$ only through
$W\cap\bigcup_{b\in\ch(v)}A_v(b)$. These quantities always refer to the
input tree, before any pruning.

\step{1 (pruning)} Fix $W$ and put $\epsilon_j=\ee^{-m_j}$. All leaves
are \emph{good}. Proceeding upward, a node $v$ with outgoing level $j$ is
good if
\begin{equation}\label{eq:good-node}
 \sum_{b\in\ch(v),\ b\text{ good}}t_{vb}(W)\ge1-\epsilon_j,
\end{equation}
and \emph{bad} otherwise. By \eqref{eq:local-masses}, at a good node the
failure mass plus the mass of bad targets is at most $\epsilon_j$. The
procedure \emph{succeeds} when the root is good.

\step{2 (potential)} Set $D_v(W)=0$ at leaves and, for $v$ with outgoing
level $j$,
\begin{equation}\label{eq:pruning-recurrence}
 D_v(W)=\epsilon_j^{-1}
 \Bigl(d_v(W)+\sum_{b\in\ch(v)}t_{vb}(W)D_b(W)\Bigr)\ge0.
\end{equation}
\emph{Claim: every bad node has $D_v(W)>1$.} Induct upward; leaves are
good. If $v$ is bad and the claim holds for its bad children, then
\[
 \epsilon_jD_v\ge d_v+\sum_{b\text{ bad}}t_{vb}D_b
 \ge d_v+\sum_{b\text{ bad}}t_{vb}
 =1-\sum_{b\text{ good}}t_{vb}>\epsilon_j,
\]
by \eqref{eq:local-masses} and the failure of \eqref{eq:good-node}.

\step{3 ($D_b$ ignores its incoming label)} Let $b$ be a child of $v$.
\emph{Claim: $D_b(W)$ depends on $W$ only through $W\setminus A_v(b)$.}
By \eqref{eq:pruning-recurrence}, $D_b$ is a function of the quantities
$d_w,t_{wc}$ for nodes $w$ in the subtree rooted at $b$. By Step~0 these
depend on $W$ only on the labels $A_w(c)$ of arcs of that subtree. Each
such arc lies on a root-to-leaf path through the arc $(v,b)$, so by
disjointness its label is disjoint from $A_v(b)$.

\step{4 (averaging over $W$)} Now let $W\sim\mu_\rho$, and let $v$ have
outgoing level $j$. By \eqref{eq:local-failure},
$\EE d_v(W)\le\ee^{-9m_j}$. For a child $b$, absence of local failure
implies $Z(Y)\ge\ee^{-10m_j}$, so
\[
 \EE\bigl[t_{vb}(W)D_b(W)\bigr]
 \le\EE\bigl[\ind_{\{\mathrm{target}=b,\ Z(Y)\ge\ee^{-10m_j}\}}D_b(W)\bigr]
 \le\ee^{11m_j}\nu_v(b)\,\EE D_b(W);
\]
the last step is \eqref{eq:distortion} with $f=D_b$, admissible by
Step~3. Inserting both bounds into \eqref{eq:pruning-recurrence}, with
$\epsilon_j^{-1}=\ee^{m_j}$,
\begin{equation}\label{eq:expected-potential}
 \EE D_v(W)\le\ee^{-8m_j}
 +\ee^{12m_j}\sum_{b\in\ch(v)}\nu_v(b)\,\EE D_b(W).
\end{equation}
This is the step that controls the dependence created by reusing $W$.

\step{5 (the root is rarely bad)} \emph{Claim:} for $v$ with outgoing
level $j$,
\begin{equation}\label{eq:unrolled}
 \EE D_v(W)\le\sum_{i=j}^k
 \exp\Bigl(12\sum_{l=j}^{i-1}m_l-8m_i\Bigr).
\end{equation}
Induct upward. For leaves (read $j=k+1$) both sides are $0$. For the
inductive step, children have outgoing level $j+1$; insert their bound
into \eqref{eq:expected-potential} and use $\sum_b\nu_v(b)=1$: the term
$\ee^{-8m_j}$ is the summand $i=j$, and $\ee^{12m_j}$ times the $i$-th
summand for $j+1$ is the $i$-th summand for $j$.

For the root $o$ ($j=1$): since $m_l\le16^{-(i-l)}m_i$, we get
$\sum_{l<i}m_l\le m_i/15$, so each exponent is at most
$(12/15-8)m_i\le-7m_i$. The $m_i$ are distinct positive integers, so by
Step~2 and Markov's inequality
\[
 \PP(o\text{ is bad})\le\PP(D_o(W)>1)\le\EE D_o(W)
 \le\sum_{m\ge1}\ee^{-7m}<\tfrac14.
\]

\step{6 (the reduced tree)} Assume the root is good. Start from the
root; at each retained non-leaf node $v$ (necessarily good), with
outgoing level $j$, let
\[
 \gamma_v(W)=\sum_{b\text{ good}}t_{vb}(W)\ge1-\epsilon_j>0.
\]
Retain exactly the good children $b$ with $t_{vb}(W)>0$, give $b$ the
weight $t_{vb}(W)/\gamma_v(W)$, and relabel the arc $(v,b)$ by
$T_b=A_v(b)\setminus W$. Checks:
\begin{itemize}[itemsep=1pt,topsep=2pt]
 \item \emph{Probability tree with $k$ levels.} Weights are positive and
 sum to $1$. Every retained non-leaf node has a retained child, since
 $\gamma_v>0$, and that child is good; so all leaves of $\TT'$ are leaves
 of $\TT$.
 \item \emph{Disjoint.} $T_b\subseteq A_v(b)$.
 \item \emph{(ii).} A path $Q$ of $\TT'$ passes through the nodes of a
 path $P$ of $\TT$, and $A_v(b)\subseteq W\cup T_b$ on each arc.
 \item \emph{Capacity.} $T_b$ is the fragment of any transition at $v$
 with target $b$. If $t_{vb}(W)>0$, some such transition is not a local
 failure, so $|T_b|\le m_j/2$, hence $|T_b|\le\lfloor m_j/2\rfloor$.
\end{itemize}

\step{7 (spread)} Fix a retained node $v$ with outgoing level $j$ and a
nonempty $J\subseteq X$. In the transition at $v$ for the fixed $W$, let
$E_v$ be the event ``no local failure and the target is good'', so
$\PP(E_v\mid W)=\gamma_v(W)$. The new child law is the law of the target
conditioned on $E_v$, and the new label is the fragment $T$. By
\eqref{eq:two-inclusions}, $T\subseteq A_v(b)$ for the target $b$, and
$b$ has law $\nu_v$ whatever $W$ is. Hence
\[
 \PP(J\subseteq T\mid E_v,W)
 =\frac{\PP(J\subseteq T,\ E_v\mid W)}{\gamma_v(W)}
 \le\frac{\PP_{b\sim\nu_v}(J\subseteq A_v(b))}{1-\epsilon_j}
 \le\frac{\alpha_j^{|J|}}{1-\epsilon_j}
 \le\Bigl(\frac{\alpha_j}{1-\epsilon_j}\Bigr)^{|J|},
\]
the last step because $|J|\ge1$. This proves (i).

The construction uses only the numbers $t_{vb}(W)$, so it is
deterministic given $W$; sources and targets serve only the analysis.
\end{proof}

\subsection{Iteration}

\begin{proof}[Proof of Theorem~\ref{thm:tree}]
\emph{Plan.} Apply Lemma~\ref{lem:reduction} repeatedly with fresh sets,
deleting levels once their labels are empty; check that the hypotheses
persist; count the attempts.

\step{1 (parameters)} If $r=1$ then $W=X$ and the claim is trivial.
Assume $r<1$ and set
\[
 a=4\sigma,\qquad\rho=\ee^{50}a,\qquad
 s=\lfloor\log_2\ell_k\rfloor+1.
\]
Since $\ell_k\ge1$, $r\ge16\ee^{50}\sigma=4\rho$, so $0<\rho<1$.

\step{2 (procedure)} Draw independent sets $W_1,W_2,\ldots$ of law
$\mu_\rho$. Keep a current tree $\TT^{(\tau)}$, where $\tau$ counts
successes and $\TT^{(0)}=\TT$. An \emph{attempt} applies
Lemma~\ref{lem:reduction} to the current tree with the next set. On
failure the tree is unchanged. On success let $\TT'$ be the output and
\emph{trim} it: capacities increase with the level, so the levels of
capacity $0$ form an initial segment; follow one downward path from the
root through these levels (chosen by a fixed rule) to a node $v^*$, and
let $\TT^{(\tau+1)}$ be the subtree of $\TT'$ rooted at $v^*$. Stop when
the current tree has no levels. Lemma~\ref{lem:reduction} itself deletes
nothing; levels disappear only by trimming. Trimming is harmless:
\begin{itemize}[itemsep=1pt,topsep=2pt]
 \item Each path $Q$ of $\TT^{(\tau+1)}$, preceded by the chosen prefix,
 is a path of $\TT'$ with the same label union, because the prefix labels
 are empty. So by \eqref{eq:recovery}, if $W$ denotes the set used,
 \begin{equation}\label{eq:trim-recovery}
  \text{for every path }Q\text{ of }\TT^{(\tau+1)}\text{ there is a path }
  P\text{ of }\TT^{(\tau)}\text{ with }\UU(P)\subseteq W\cup\UU(Q).
 \end{equation}
 \item Nodes of $\TT^{(\tau+1)}$ keep their children, child laws and
 labels from $\TT'$. So the tree is a disjoint probability tree, and the
 capacity and spread bounds of Lemma~\ref{lem:reduction}(i) persist on
 the remaining levels.
\end{itemize}

\step{3 (invariants)} Index levels by their index $i$ in $\TT$. After
$\tau$ successes, by induction and
$\lfloor\lfloor y\rfloor/2\rfloor=\lfloor y/2\rfloor$:
\begin{itemize}[itemsep=1pt,topsep=2pt]
 \item level $i$ remains if and only if
 $m_i(\tau)=\lfloor\ell_i/2^\tau\rfloor\ge1$, and then has capacity
 $m_i(\tau)$;
 \item remaining capacities satisfy
 $m_{i+1}(\tau)\ge\lfloor16\ell_i/2^\tau\rfloor\ge16m_i(\tau)$;
 \item a remaining level $i$ is $\alpha_i(\tau)$-spread, where
 $\alpha_i(\tau)=\sigma\prod_{u=0}^{\tau-1}(1-\ee^{-m_i(u)})^{-1}$, by
 Lemma~\ref{lem:reduction}(i) applied at each earlier success.
\end{itemize}

\step{4 ($\alpha_i(\tau)\le a$)} For a remaining level,
$m_i(0)>m_i(1)>\cdots>m_i(\tau-1)\ge1$, because
$\lfloor y/2\rfloor<y$ for integers $y\ge1$. So the factors of
$\alpha_i(\tau)$ belong to distinct positive integers $m$, and
$\alpha_i(\tau)\le\sigma\prod_{m\ge1}(1-\ee^{-m})^{-1}$. Using
$-\log(1-z)\le z/(1-z)\le z/(1-\ee^{-1})$ for $0<z\le\ee^{-1}$,
\[
 \sum_{m\ge1}-\log(1-\ee^{-m})
 \le\frac1{1-\ee^{-1}}\sum_{m\ge1}\ee^{-m}
 =\frac{\ee^{-1}}{(1-\ee^{-1})^2}<0.93<\log4,
\]
so the product is below $4$ and $\alpha_i(\tau)\le4\sigma=a$. By Steps
1--4, as long as it has a level, the current tree satisfies the
hypotheses of Lemma~\ref{lem:reduction} with these $a,\rho$. The current
tree is determined by the earlier sets and the next set is independent of
them; so each attempt succeeds with conditional probability at least
$3/4$ given the past.

\step{5 (termination and recovery)} Level $k$ remains exactly while
$2^\tau\le\ell_k$, that is $\tau\le s-1$. So the procedure stops after
exactly $s$ successes, and $\TT^{(s)}$ is a single node whose only path
$Q_s$ is empty. Let $W^{(1)},\ldots,W^{(s)}$ be the sets used at the
successes. Applying \eqref{eq:trim-recovery} for $\tau=s,s-1,\ldots,1$
gives paths $Q_{\tau-1}$ of $\TT^{(\tau-1)}$ with
$\UU(Q_{\tau-1})\subseteq W^{(\tau)}\cup\UU(Q_\tau)$. Hence
\[
 \UU(Q_0)\subseteq W^{(1)}\cup\cdots\cup W^{(s)},
 \qquad Q_0\text{ a path of }\TT.
\]

\step{6 (number of attempts)} Let $G_\tau$ be the number of attempts
after the $(\tau-1)$-st success up to and including the $\tau$-th. The
event $G_\tau>j$ requires $j$ consecutive failures, each of conditional
probability at most $1/4$ given all earlier sets, so
$\PP(G_\tau>j)\le4^{-j}$ and
$\EE G_\tau=\sum_{j\ge0}\PP(G_\tau>j)\le4/3$. The total number of
attempts $N=G_1+\cdots+G_s$ has $\EE N\le4s/3$, and by Markov's
inequality $\PP(N>4s)\le1/3$.

\step{7 (density)} Let $W_*=W_1\cup\cdots\cup W_{4s}$. It has law
$\mu_{r_*}$ with
\[
 r_*=1-(1-\rho)^{4s}\le4s\rho=16\ee^{50}\sigma s
 \le16\ee^{50}\sigma(1+\log_2\ell_k)=r.
\]
On $\{N\le4s\}$, of probability at least $2/3$, all successful sets are
among $W_1,\ldots,W_{4s}$, so $\UU(Q_0)\subseteq W_*$ by Step~5. The
event that a set contains some path union is closed under supersets, so
by monotonicity the same bound holds for $\mu_r$.
\end{proof}

\section{Application to graph containment}\label{sec:graphs}

The construction below is the normalized-containment maximization and
conditional-extension argument of
Tran~\cite[Section~3, Lemma~3.1]{Tran26}, with constants suited to the
factor $16$; identity \eqref{eq:copy-ratio} also appears
in~\cite[Equation~(5)]{MNSZ22}.

\paragraph{Conventions.}
Put $q=\pE(n,H)$, $\sigma=128q$ and $X=\binom{[n]}2$. In this section a
graph without isolated vertices is identified with its edge set; $H$ now
denotes the edge set of the given graph, and $F\subseteq H$ means an edge
subset. A \emph{copy} of such an $F$ is a subset of $X$ isomorphic to
$F$. Let $M(F)=N(K_n,F)$ be the number of copies of $F$, and
$M(I,F)=N(I,F)$ the number of subsets of $I$ isomorphic to $F$, with
$M(\varnothing)=M(I,\varnothing)=1$. The number $q$ always refers to the
original graph. Two remarks justify this:
\begin{itemize}[itemsep=1pt,topsep=2pt]
 \item If an edge set on $[n]$ contains a copy of the edge set $H$, the
 corresponding graph contains the original graph: place its isolated
 vertices on unused vertices, which exist since the original graph has at
 most $n$ vertices, and copies need not be induced.
 \item Each nonempty $F\subseteq H$, on its incident vertices, is a
 subgraph of the original graph. Since the infimum defining $q$ is
 attained, fact~\ref{fact:a} gives
 \begin{equation}\label{eq:first-moment}
  M(F)\,q^{|F|}\ge\tfrac12.
 \end{equation}
\end{itemize}

\paragraph{Copy-ratio identity.}
Let $F\subseteq I\subseteq H$, let $\boldsymbol I$ be a uniformly random
copy of $I$, and let $F_0$ be a fixed copy of $F$. Then
\begin{equation}\label{eq:copy-ratio}
 \PP(F_0\subseteq\boldsymbol I)=\frac{M(I,F)}{M(F)}.
\end{equation}
Indeed, permutations of $[n]$ preserve the law of $\boldsymbol I$ and act
transitively on the copies of $F$, so the left side is the same number
$\pi$ for every $F_0$. Every value of $\boldsymbol I$ contains exactly
$M(I,F)$ copies of $F$; taking expectations in
$\sum_{F_0}\ind_{\{F_0\subseteq\boldsymbol I\}}=M(I,F)$ gives
$M(F)\pi=M(I,F)$. In particular the probability depends only on the
isomorphism types of $F$ and $I$.

\begin{lemma}[Graph hierarchy]\label{lem:hierarchy}
There is a disjoint probability tree on $X=\binom{[n]}2$ whose path
unions are copies of $H$ with its isolated vertices removed, such that
every level is $\sigma$-spread and its label sizes
$\ell_1,\ldots,\ell_k$ satisfy
\[
 1\le\ell_1,\qquad\ell_{i+1}\ge16\ell_i,\qquad\ell_k\le e(H).
\]
\end{lemma}

This is the input to Theorem~\ref{thm:tree} in the proof of
Theorem~\ref{thm:main}.

\begin{proof}
\emph{Plan.} Choose each $H_{i-1}\subseteq H_i$ by a maximization that
forces $|H_{i-1}|<|H_i|/17$ (from first moments) and conditional spread
(from maximality).

\step{1 (maximizers are small)} For nonempty $I\subseteq H$ and
$S\subseteq I$ let
\begin{equation}\label{eq:maximizer}
 R_\sigma(I,S)=\sigma^{-|S|}\,\PP(S_0\subseteq\boldsymbol I),
\end{equation}
where $S_0$ is any copy of $S$ and $\boldsymbol I$ a uniform copy of $I$;
by \eqref{eq:copy-ratio} this does not depend on $S_0$. Note
$R_\sigma(I,\varnothing)=1$. \emph{Claim: if $x=|S|\ge|I|/17$, then
$R_\sigma(I,S)<1$.} Here $x\ge1$. By \eqref{eq:copy-ratio}, by
$1/M(S)\le2q^x$ from \eqref{eq:first-moment}, and by
$M(I,S)\le\binom{|I|}{x}\le(\ee|I|/x)^x\le(17\ee)^x$,
\[
 R_\sigma(I,S)\le\sigma^{-x}\cdot2q^x(17\ee)^x
 =2\Bigl(\frac{17\ee}{128}\Bigr)^x\le\frac{34\ee}{128}<0.73.
\]
Hence every maximizer $S$ of $R_\sigma(I,\cdot)$, whose value is at least
$1$, satisfies $|S|<|I|/17$.

\step{2 (chain)} Starting from $I=H$, repeatedly replace the current
nonempty $I$ by a maximizer of $R_\sigma(I,\cdot)$ (ties broken by a
fixed rule) until $\varnothing$ is reached. In reverse order this gives
\[
 \varnothing=H_0\subset H_1\subset\cdots\subset H_k=H,\qquad
 H_{i-1}\text{ maximizes }R_\sigma(H_i,\cdot),\qquad
 |H_{i-1}|<|H_i|/17,
\]
with $k\ge1$.

\step{3 (tree)} The root is the empty edge set. A node at depth $i-1$ is
a copy $S_0$ of $H_{i-1}$ together with its history. Its children are all
copies $I_0$ of $H_i$ with $S_0\subseteq I_0$, with uniform law, and the
arc is labelled $I_0\setminus S_0$. Children exist: take an isomorphism
from $H_{i-1}$ to $S_0$ and extend it injectively to the vertices of
$H_i$, which is possible because $H$ has at most $n$ vertices; the image
of $H_i$ is a copy containing $S_0$. Along a path the copies are nested,
so the labels are pairwise disjoint and their union is a copy of $H$. All
leaves have depth $k$.

\step{4 (sizes)} Every label of level $i$ has size
$\ell_i=|H_i|-|H_{i-1}|\ge1$. Moreover $\ell_k\le|H|=e(H)$ and
$\ell_{i+1}>17|H_i|-|H_i|=16|H_i|\ge16\ell_i$.

\step{5 (spread)} Fix a node $S_0$, a copy of $H_{i-1}$, let $A$ be its
random label, and let $\boldsymbol I$ be a uniform copy of $H_i$.
Conditioned on $S_0\subseteq\boldsymbol I$, $\boldsymbol I$ is uniform on
the children of $S_0$. Let $J\subseteq X$ be nonempty. If
$J\cap S_0\ne\varnothing$ then $\PP(J\subseteq A)=0$. Otherwise
\begin{equation}\label{eq:conditional-ratio}
 \PP(J\subseteq A)=
 \frac{\PP(S_0\cup J\subseteq\boldsymbol I)}
      {\PP(S_0\subseteq\boldsymbol I)}.
\end{equation}
If the numerator is $0$ we are done. Otherwise $S_0\cup J$ lies in some
copy $I_0$ of $H_i$; an isomorphism $I_0\to H_i$ maps $S_0\cup J$ onto an
edge subset $U\subseteq H_i$ with $|U|=|S_0|+|J|$, and $S_0\cup J$ is a
copy of $U$. By \eqref{eq:maximizer}, the numerator is
$\sigma^{|S_0|+|J|}R_\sigma(H_i,U)$ and the denominator is
$\sigma^{|S_0|}R_\sigma(H_i,H_{i-1})$. Maximality of $H_{i-1}$ gives
$R_\sigma(H_i,U)\le R_\sigma(H_i,H_{i-1})$, hence
$\PP(J\subseteq A)\le\sigma^{|J|}$.
\end{proof}

\emph{Sanity check.} If $H$ is a single edge, then $q=1/(n(n-1))$, $k=1$,
and the root has $\binom n2$ equally likely children with one-edge
labels; indeed $1/\binom n2=2q\le\sigma$.

\begin{proof}[Proof of Theorem~\ref{thm:main}]
Apply Theorem~\ref{thm:tree} to the tree of Lemma~\ref{lem:hierarchy},
with $\sigma=128q$. Since $\ell_k\le h$, its density
$\min\{1,16\ee^{50}\cdot128q(1+\log_2\ell_k)\}$ is at most
\[
 r=\min\{1,\;2048\ee^{50}q(1+\log_2h)\}.
\]
The edge set of $G(n,r)$ has law $\mu_r$ on $X$. By
Theorem~\ref{thm:tree} and monotonicity, with probability at least $2/3$
it contains a path union, which is a copy of the edge set $H$; by the
first remark of the Conventions, $G(n,r)$ then contains the original
graph. As $2/3\ge1/2$, $\pc(n,H)\le r$, which is \eqref{eq:stronger}.
Finally $h\le\binom n2\le n^2$ and $\log_2n\ge1$ give
$1+\log_2h\le1+2\log_2n\le3\log_2n$, and $3\cdot2048=6144$.
\end{proof}

\section*{Verification checklist}

No step is marked unverified; the context paragraphs of
Section~\ref{sec:introduction} were not checked against the cited
papers. A sceptical reader should check these first:
\begin{enumerate}[itemsep=2pt,topsep=3pt]
 \item \textbf{Cancellation and independence in Lemma~\ref{lem:local}}
 (Steps 3, 4, 6): $\rho^{-|A_b|}\rho^{|A_b\setminus A_c|}
 =\rho^{-|A_c\cap A_b|}$, and the indicator of
 $A_b\setminus A_c\subseteq W$ is independent of any function of
 $W\setminus A_b$.
 \item \textbf{$D_b$ ignores $A_v(b)$} (Lemma~\ref{lem:reduction},
 Step~3). This is the only place where disjointness along paths enters
 the probability estimates; it licenses \eqref{eq:distortion} with
 $f=D_b$ although one $W$ serves all nodes.
 \item \textbf{Exponent bookkeeping} (Lemma~\ref{lem:reduction}, Steps
 4--5): $\ee^{m}\cdot\ee^{-9m}$ and $\ee^{m}\cdot\ee^{11m}$ in
 \eqref{eq:expected-potential}, then $12/15-8\le-7$ and
 $\sum_{m\ge1}\ee^{-7m}<1/4$.
 \item \textbf{Spread after conditioning} (Lemma~\ref{lem:reduction},
 Step~7): it rests on $T\subseteq A_v(b)$ and on the target having law
 $\nu_v$ for every fixed $W$; the cost is only the factor
 $1/\gamma_v\le1/(1-\ee^{-m_j})$.
 \item \textbf{Hypotheses persist under iteration}
 (Theorem~\ref{thm:tree}, Steps 2--4): trimming keeps
 \eqref{eq:trim-recovery}, the ratio $16$ survives the floors, and the
 accumulated spread factor is below $4$ because the capacities at one
 level are distinct integers.
 \item \textbf{Hierarchy} (Lemma~\ref{lem:hierarchy}, Steps 1 and 5):
 $2\cdot17\ee/128<1$ forces $|H_{i-1}|<|H_i|/17$; the spread bound uses
 that \eqref{eq:copy-ratio} depends only on isomorphism types, so
 maximality applies to $U\cong S_0\cup J$.
 \item \textbf{Constants}: $\rho=4\ee^{50}\sigma$, at most $4s$ sets,
 $\sigma=128q$, giving $4\cdot4\cdot128=2048$; then $3\cdot2048=6144$.
\end{enumerate}

\begin{thebibliography}{99}
\raggedright

\bibitem{ALWZ21}
R.~Alweiss, S.~Lovett, K.~Wu, and J.~Zhang,
\emph{Improved bounds for the sunflower lemma},
Annals of Mathematics (2) \textbf{194} (2021), no.~3, 795--815.
\href{https://doi.org/10.4007/annals.2021.194.3.5}{doi:10.4007/annals.2021.194.3.5}.

\bibitem{Bollobas81}
B.~Bollob\'as,
\emph{Threshold functions for small subgraphs},
Mathematical Proceedings of the Cambridge Philosophical Society
\textbf{90} (1981), no.~2, 197--206.
\href{https://doi.org/10.1017/S0305004100058655}{doi:10.1017/S0305004100058655}.

\bibitem{DKP25}
Q.~Dubroff, J.~Kahn, and J.~Park,
\emph{On the ``second'' Kahn--Kalai conjecture},
arXiv:2508.14269v1 (August 19, 2025).
\url{https://arxiv.org/abs/2508.14269v1}.

\bibitem{DKP25Cases}
Q.~Dubroff, J.~Kahn, and J.~Park,
\emph{On the ``second'' Kahn--Kalai conjecture: cliques, cycles, and trees},
arXiv:2510.02005v2 (October 23, 2025).
\url{https://arxiv.org/abs/2510.02005v2}.

\bibitem{ER60}
P.~Erd\H{o}s and A.~R\'enyi,
\emph{On the evolution of random graphs},
Publications of the Mathematical Institute of the Hungarian Academy of
Sciences \textbf{5} (1960), 17--61.
\url{https://static.renyi.hu/~p_erdos/1960-10.pdf}.

\bibitem{ER66}
P.~Erd\H{o}s and A.~R\'enyi,
\emph{On the existence of a factor of degree one of a connected random graph},
Acta Mathematica Academiae Scientiarum Hungaricae
\textbf{17} (1966), 359--368.
\url{https://users.renyi.hu/~p_erdos/1966-16.pdf}.

\bibitem{FKNP21}
K.~Frankston, J.~Kahn, B.~Narayanan, and J.~Park,
\emph{Thresholds versus fractional expectation-thresholds},
Annals of Mathematics (2) \textbf{194} (2021), no.~2, 475--495.
\href{https://doi.org/10.4007/annals.2021.194.2.2}{doi:10.4007/annals.2021.194.2.2}.

\bibitem{KK07}
J.~Kahn and G.~Kalai,
\emph{Thresholds and expectation thresholds},
Combinatorics, Probability and Computing
\textbf{16} (2007), no.~3, 495--502.
\href{https://doi.org/10.1017/S0963548307008474}{doi:10.1017/S0963548307008474}.

\bibitem{LSZ19}
S.~Lovett, N.~Solomon, and J.~Zhang,
\emph{From DNF compression to sunflower theorems via regularity},
34th Computational Complexity Conference, LIPIcs \textbf{137} (2019),
5:1--5:14.
\href{https://doi.org/10.4230/LIPIcs.CCC.2019.5}{doi:10.4230/LIPIcs.CCC.2019.5}.

\bibitem{MNSZ25}
E.~Mossel, J.~Niles-Weed, N.~Sun, and I.~Zadik,
\emph{A Bayesian proof of the spread lemma},
Random Structures \& Algorithms \textbf{66} (2025), no.~4, e70008.
\href{https://doi.org/10.1002/rsa.70008}{doi:10.1002/rsa.70008}.

\bibitem{MNSZ22}
E.~Mossel, J.~Niles-Weed, N.~Sun, and I.~Zadik,
\emph{On the second Kahn--Kalai conjecture},
arXiv:2209.03326v1 (September 7, 2022).
\url{https://arxiv.org/abs/2209.03326v1}.

\bibitem{OpenAIThresholdEquivalence2026}
OpenAI,
\emph{Integral and fractional expectation thresholds are equivalent},
OpenAI Math Release preprint
\href{https://github.com/openai/math/blob/main/preprints/Integral-and-fractional-expectation-thresholds-are-equivalent-September-23-2026/paper.pdf}{OAI:Integral-and-fractional-expectation-thresholds-are-equivalent-September-23-2026},
2026, Theorem~1.1, p.~2.

\bibitem{PP24}
J.~Park and H.~T.~Pham,
\emph{A proof of the Kahn--Kalai conjecture},
Journal of the American Mathematical Society
\textbf{37} (2024), no.~1, 235--243.
\href{https://doi.org/10.1090/jams/1028}{doi:10.1090/jams/1028}.

\bibitem{Riordan00}
O.~Riordan,
\emph{Spanning subgraphs of random graphs},
Combinatorics, Probability and Computing
\textbf{9} (2000), no.~2, 125--148.
\href{https://doi.org/10.1017/S0963548399004150}{doi:10.1017/S0963548399004150}.

\bibitem{Tal10}
M.~Talagrand,
\emph{Are many small sets explicitly small?},
Proceedings of the 42nd ACM Symposium on Theory of Computing (2010),
13--36.
\href{https://doi.org/10.1145/1806689.1806693}{doi:10.1145/1806689.1806693}.

\bibitem{Tran26}
T.~Tran,
\emph{Fractional expectation thresholds and the ``second'' Kahn--Kalai
conjecture},
arXiv:2609.20546v2 (September 24, 2026).
\url{https://arxiv.org/abs/2609.20546v2}.

\end{thebibliography}

\end{document}
